%package list
\documentclass{article}
\usepackage[top=3cm, bottom=3cm, outer=3cm, inner=3cm]{geometry}
\usepackage{graphicx}
\usepackage{url}
\usepackage{hyperref}
\usepackage{array}
\newcolumntype{x}[1]{>{\centering\arraybackslash\hspace{0pt}}p{#1}}
\usepackage{natbib}
\usepackage{pdfpages}
\usepackage{multirow}
\usepackage[T1]{fontenc}
\usepackage{imakeidx}
\usepackage{wrapfig}

\usepackage{hyperref}
\hypersetup{
    colorlinks=true,
    linkcolor=black,
    filecolor=magenta,
    urlcolor=blue,
    pdftitle={Pruebas Mockito},
    pdfpagemode=FullScreen,
    }

\urlstyle{same}

\usepackage[normalem]{ulem}
\useunder{\uline}{\ul}{}

\usepackage{minted}

\usepackage{listings}
\usepackage{color, colortbl}
\definecolor{dkgreen}{rgb}{0,0.6,0}
\definecolor{gray}{rgb}{0.5,0.5,0.5}
\definecolor{mauve}{rgb}{0.58,0,0.82}
\definecolor{codebackground}{rgb}{0.95, 0.95, 0.92}
\definecolor{tablebackground}{rgb}{0.0, 0.45, 0.63}
\lstset{frame=tb,
    language=bash,
    aboveskip=3mm,
    belowskip=3mm,
    showstringspaces=false,
    columns=flexible,
    basicstyle={\small\ttfamily},
    numbers=none,
    numberstyle=\tiny\color{gray},
    keywordstyle=\color{blue},
    commentstyle=\color{dkgreen},
    stringstyle=\color{mauve},
    breaklines=true,
    breakatwhitespace=true,
    tabsize=3,
    backgroundcolor= \color{codebackground},
}

%%%%%%%%%%%%%%%%%%%%%%%%%%%%%%%%%%%%%%%%%%%%%%%%%%%%%%%%%%%%%%%%%%%%%%%%%%%%
\newcommand{\csemail}{pramirezsa@ulasalle.edu.pe}
\newcommand{\csdocente}{MSc. Luis Pati\~no Herrera}
\newcommand{\cscurso}{Pruebas de Software}
\newcommand{\csuniversidad}{Universidad La Salle}
\newcommand{\csescuela}{Escuela Profesional de Ingenier\'ia de Software}
\newcommand{\cspracnr}{02}
\newcommand{\cstema}{Pruebas con Mockito}
%%%%%%%%%%%%%%%%%%%%%%%%%%%%%%%%%%%%%%%%%%%%%%%%%%%%%%%%%%%%%%%%%%%%%%%%%%%%

\usepackage[english,spanish]{babel}
\usepackage[utf8]{inputenc}
\AtBeginDocument{\selectlanguage{spanish}}
\renewcommand{\figurename}{Figura}
\renewcommand{\refname}{Referencias}
\renewcommand{\tablename}{Tabla}
\AtBeginDocument{%
    \renewcommand\tablename{Tabla}
}

\usepackage{fancyhdr}
\pagestyle{fancy}
\fancyhf{}
\setlength{\headheight}{30pt}
\renewcommand{\headrulewidth}{1pt}
\renewcommand{\footrulewidth}{1pt}
\fancyhead[L]{\raisebox{-0.2\height}{\includegraphics[width=3cm]{logo_ulasalle (1).png}}}
\fancyhead[C]{}
\fancyhead[R]{\fontsize{7}{7}\selectfont    \csuniversidad \\ \csescuela \\ \textbf{\cscurso} }
\fancyfoot[L]{}
\fancyfoot[C]{Pruebas de Software}
\fancyfoot[R]{P\'agina \thepage}

\begin{document}

    \vspace*{10px}

    \begin{center}
        \fontsize{17}{17} \textbf{ Pr\'actica \cspracnr}
    \end{center}

\renewcommand{\arraystretch}{1.5}
\begin{table}[h]
    \begin{tabular}{|x{4.7cm}|x{4.8cm}|x{4.8cm}|}
        \hline
        \textbf{DOCENTE} & \textbf{CARRERA}  & \textbf{CURSO}   \\
        \hline
        \csdocente & \csescuela & \cscurso    \\
        \hline
    \end{tabular}
\end{table}

\begin{table}[h]
    \begin{tabular}{|x{4.7cm}|x{4.8cm}|x{4.8cm}|}
        \hline
        \textbf{GRUPO} & \textbf{TEMA}  & \textbf{DURACI\'ON}   \\
        \hline
        \ - & \cstema & -  \\
        \hline
    \end{tabular}
\end{table}
\renewcommand{\arraystretch}{1}

    \section*{Integrantes}
    \begin{itemize}
        \item Patrick Andres Ramirez Santos
    \end{itemize}

    \section*{Repositorio}
    C\'odigo fuente disponible en: \url{URL\_DEL\_REPOSITORIO\_GITHUB}

    \tableofcontents

\newpage

%% ============================================================
\section{Objetivo}
Implementar el tutorial de Mockito de Vogella utilizando JUnit 5, demostrando el uso de mocks, stubs, verificaciones e inyecci\'on de dependencias en pruebas unitarias.

Tutorial base: \url{https://www.vogella.com/tutorials/Mockito/article.html}

%% ============================================================
\section{Estructura del proyecto}

El proyecto sigue la estructura est\'andar de Maven. Las clases fuente est\'an en \texttt{src/main/java/org/example/} y las pruebas en \texttt{src/test/java/org/example/}. Las dependencias principales son \texttt{mockito-core}, \texttt{mockito-junit-jupiter} y \texttt{junit-jupiter}.

\begin{center}
    \textit{Estructura del proyecto en el IDE o terminal}

    %\includegraphics[width=0.8\linewidth]{PLACEHOLDER_ESTRUCTURA.png}
    \textbf{[CAPTURA: estructura de carpetas del proyecto]}
\end{center}

%% ============================================================
\section{Compilaci\'on}

Se verifica que todo el c\'odigo fuente compile sin errores:

\begin{minted}{bash}
$ mvn compile
\end{minted}

\begin{center}
    \textit{Compilaci\'on exitosa del proyecto}

    %\includegraphics[width=0.9\linewidth]{PLACEHOLDER_COMPILE.png}
    \textbf{[CAPTURA: mvn compile -- BUILD SUCCESS]}
\end{center}

%% ============================================================
\section{Ejecuci\'on de pruebas}

\subsection{Todas las pruebas}

Se ejecutan las 6 pruebas del proyecto con:

\begin{minted}{bash}
$ mvn test
\end{minted}

\begin{center}
    \textit{Ejecuci\'on de todas las pruebas (6 tests, 0 fallos)}

    %\includegraphics[width=0.9\linewidth]{PLACEHOLDER_TEST_ALL.png}
    \textbf{[CAPTURA: mvn test -- Tests run: 6, BUILD SUCCESS]}
\end{center}

\subsection{UserServiceTest -- Mocks y verificaciones}

Prueba los conceptos b\'asicos: \texttt{@Mock}, \texttt{when/thenReturn} y \texttt{verify}. Se mockean \texttt{UserRepository} y \texttt{EmailService} para probar \texttt{registerUser()} sin dependencias reales.

\begin{minted}{bash}
$ mvn test -Dtest=UserServiceTest
\end{minted}

\begin{center}
    \textit{UserServiceTest: mockeo de interfaces UserRepository y EmailService}

    %\includegraphics[width=0.9\linewidth]{PLACEHOLDER_TEST_USER.png}
    \textbf{[CAPTURA: mvn test -Dtest=UserServiceTest -- 2 tests passed]}
\end{center}

\subsection{VolumeUtilTests -- Mock inline y verifyNoMoreInteractions}

Demuestra la creaci\'on de mocks con \texttt{mock()} (en l\'inea) y la verificaci\'on estricta con \texttt{verifyNoMoreInteractions} para asegurar que en modo silencioso no se modifica el volumen.

\begin{minted}{bash}
$ mvn test -Dtest=VolumeUtilTests
\end{minted}

\begin{center}
    \textit{VolumeUtilTests: mock de AudioManager con verificaci\'on estricta}

    %\includegraphics[width=0.9\linewidth]{PLACEHOLDER_TEST_VOLUME.png}
    \textbf{[CAPTURA: mvn test -Dtest=VolumeUtilTests -- 2 tests passed]}
\end{center}

\subsection{ConfigureThreadingUtilTests -- verify con only()}

Verifica que \texttt{configureThreadPool()} llame \'unicamente a \texttt{getNumberOfThreads()} usando \texttt{verify(app, only())}.

\begin{minted}{bash}
$ mvn test -Dtest=ConfigureThreadingUtilTests
\end{minted}

\begin{center}
    \textit{ConfigureThreadingUtilTests: verificaci\'on con only() sobre MyApplication}

    %\includegraphics[width=0.9\linewidth]{PLACEHOLDER_TEST_THREADING.png}
    \textbf{[CAPTURA: mvn test -Dtest=ConfigureThreadingUtilTests -- 1 test passed]}
\end{center}

\subsection{ArticleManagerTest -- @InjectMocks}

Demuestra inyecci\'on autom\'atica de mocks con \texttt{@InjectMocks}. Mockito inyecta \texttt{ArticleDatabase} y \texttt{User} en el constructor de \texttt{ArticleManager}.

\begin{minted}{bash}
$ mvn test -Dtest=ArticleManagerTest
\end{minted}

\begin{center}
    \textit{ArticleManagerTest: inyecci\'on autom\'atica con @InjectMocks en ArticleManager}

    %\includegraphics[width=0.9\linewidth]{PLACEHOLDER_TEST_ARTICLE.png}
    \textbf{[CAPTURA: mvn test -Dtest=ArticleManagerTest -- 1 test passed]}
\end{center}

%% ============================================================
\section{Conclusiones}
\begin{itemize}
    \item Se implement\'o el tutorial completo de Mockito con 4 clases de prueba y 6 tests.
    \item \texttt{@Mock} y \texttt{when/thenReturn} permiten simular dependencias sin implementaciones reales.
    \item \texttt{verify} y \texttt{verifyNoMoreInteractions} aseguran que los mocks se usen correctamente.
    \item \texttt{@InjectMocks} simplifica la inyecci\'on de dependencias en las pruebas.
\end{itemize}

\end{document}
